\section{Limitations and conclusion}
\label{sec:conclusion}

The exact results concern linear spectral analysis before a nonlinear readout. Finite Chebyshev order introduces frame and reconstruction error, particularly when a mapped pole approaches the spectral interval. Finite-spectrum product-sum interpolation may require one term per distinct eigenvalue and therefore does not establish parameter efficiency. The learned complex lift is a feature embedding rather than a graph analytic signal. Finally, norm identities do not establish resistance to oversquashing, stable nonlinear optimization, or long-range reasoning.

The H100 benchmark rows are preliminary diagnostics, not confirmatory comparisons. Each entry is one legacy run on one split; the local implementations were not checked against their upstream repositories, no run manifest was saved, and the strict artifact audit found two metric drifts. AUROC values were therefore recomputed from saved predictions using a tie-aware one-vs-rest estimator. These results provide no uncertainty estimate, significance test, or evidence that the relaxed GBDN+ architecture outperforms established methods. The matched multi-seed benchmark needed for such a claim remains future work. Compute reporting is also incomplete outside the two Peptides-func runs. Because the work is foundational filter-bank research without a deployment target, we do not identify a direct societal impact; application-specific fairness, privacy, and safety properties would require separate evaluation.

Within this scope, learned Blaschke--Cayley factors provide an interpretable all-pass parameterization for a reconstructing graph filter bank. The theory connects root geometry to phase, mapped target poles, finite-order error, complementary energy separation, and multilevel synthesis while explicitly separating complete-map conditioning from target-specific information transmission. The current controlled studies illustrate the proposed mechanisms but require provenance-complete reruns before supporting contribution claims; the legacy downstream runs establish only a preliminary portability snapshot.
